\documentclass[11pt]{article}
\usepackage[margin=1in]{geometry}
\usepackage{amsmath,amssymb,amsthm}
\newtheorem{proposition}{Proposition}
\newcommand{\vol}{\operatorname{vol}}
\title{One-edge multiway repair can require a remote split}
\author{Working research note; no historical priority claim}
\date{November 15, 2026}
\begin{document}
\maketitle
For a nonempty vertex set $S$, use the one-sided conductance
$\phi_G(S)=w(S,V\setminus S)/\vol_G(S)$, even when $S$ contains more
than half of the volume. This note records an explicit obstruction to a
particularly tempting locality claim: retaining all old clusters not
incident to an edited edge and repairing only the endpoint clusters.
It is not a lower bound for update time in an implicit data structure.

\begin{proposition}[An eight-vertex counterexample]
For every $0<\eta<1$, let $G_\eta$ be the weighted path whose vertex
order is
\[
 u,a,s,r,q,p,b,v
\]
and whose successive edge weights are
\[
 1,\eta,1,\eta,1,\eta,1.
\]
The partition
\[
 C_1=\{u,a\},\qquad C_2=\{b,v\},\qquad
 C_3=\{p,q,r,s\}
\]
has maximum conductance $\eta/(2+\eta)$.
Insert the unit-weight edge $uv$ to form $G'_\eta$, and put
$U=C_1\cup C_2$. Then:
\begin{enumerate}
\item $G'_\eta$ has a three-way partition of maximum conductance
at most $\eta/(1+\eta)$, and its normalized Laplacian satisfies
$\lambda_3(G'_\eta)\le2\eta/(1+\eta)$;
\item every pair of disjoint nonempty subsets of $U$ contains a set
of conductance at least $1/3$ in $G'_\eta$.
\end{enumerate}
Consequently, no repair procedure retaining $C_3$ unchanged and
partitioning only $U$ into two new clusters can guarantee maximum
conductance tending to zero with $\eta$, although $\lambda_3(G'_\eta)$
tends to zero and a suitable global three-way partition exists.
\end{proposition}
\begin{proof}
In the old graph, the degrees give
$\phi(C_1)=\phi(C_2)=\eta/(2+\eta)$ and
$\phi(C_3)=2\eta/(4+4\eta)\le\eta/(2+\eta)$.

In the new graph use the partition
\[
 U,\qquad A=\{p,q\},\qquad B=\{r,s\}.
\]
Its conductances are respectively
$\eta/(3+\eta)$, $\eta/(1+\eta)$, and $\eta/(1+\eta)$.
For completeness, the elementary variational implication
$\lambda_k\le2\max_i\phi(S_i)$ follows by taking the span of the
indicators of disjoint $S_i$: for coefficients $c_i$, the energy across
each interset edge is at most $2w(c_i^2+c_j^2)$, and the mass is
$\sum_i c_i^2\vol(S_i)$. This proves the stated eigenvalue bound.

Of any two disjoint nonempty subsets of $U$, at least one omits $u$ or
$v$; call it $S$. If it contains neither $u$ nor $v$, it is a nonempty
subset of $\{a,b\}$, and $\phi(S)=1$.
Otherwise, by symmetry assume it contains $u$ and omits $v$. The four
possibilities and conductances are
\[
 \begin{array}{c|c}
 S&\phi_{G'_\eta}(S)\\\hline
 \{u\}&1\\
 \{u,a\}&(1+\eta)/(3+\eta)\\
 \{u,b\}&1\\
 \{u,a,b\}&(2+2\eta)/(4+2\eta).
 \end{array}
\]
All are at least $1/3$.
\end{proof}

\paragraph{Leverage scope.}
The edit in this example has leverage
$\tau=R_{uv}(G_\eta)=4+3/\eta$, which diverges as $\eta\downarrow0$.
It therefore does not contradict the useful elementary stability
principle below. A fixed-leverage obstruction of comparable strength
for constant-factor approximate conductance would be impossible.

\begin{proposition}[What bounded leverage already gives for free]
Let $G'$ be obtained from a connected weighted graph $G$ by adding
weight $w>0$ on $uv$, and put $\tau=wR_{uv}(G)$. For every nonempty
$S$, one has
\[
 \frac{\phi_G(S)}{1+\tau}\le\phi_{G'}(S)
 \le(1+\tau)\phi_G(S).
\]
Thus an $a$-approximate solution to any fixed family of feasible
multiway conductance partitions in $G$ remains an
$a(1+\tau)^2$-approximate solution in $G'$, without relabeling vertices.
The same assertion applies when feasibility uses fixed vertex masses;
changing-volume balance constraints need separate treatment.
\end{proposition}
\begin{proof}
On the constants quotient,
$wbb^T\preceq\tau L$, because
$(b^Tx)^2\le R_{uv}x^TLx$. Hence
$L\preceq L'\preceq(1+\tau)L$.
Also $R_{uv}\ge1/d_u,1/d_v$, so
$D\preceq D'\preceq(1+\tau)D$.
Apply the two comparisons to the numerator and denominator of
$\phi(S)=\mathbf1_S^TL\mathbf1_S/(\mathbf1_S^TD\mathbf1_S)$.
Taking maxima and then minima over the same feasible family proves
the approximation statement.
\end{proof}

\paragraph{Implication for a research target.}
A meaningful dynamic theorem should improve on this automatic transfer
by obtaining a $(1+\varepsilon)$ approximation to the desired objective,
or by controlling many-update accumulation, query cost, or recourse
under an explicit structural assumption. Constant-factor reuse after
one bounded-leverage insertion is already a direct comparison argument.
The eight-vertex obstruction shows that endpoint-only repair is too
strong without such an assumption.
\end{document}
